\documentclass[a4paper,12pt]{article}
%need to format subitems

\setcounter{tocdepth}{2} % show only 2 levels in toc

\begin{document}

\tableofcontents
\newpage

\section{Capital}

\subsection{Chapter 13: Co-operation}
\textit{my thoughts on the chapter}

\begin{itemize}
\item
Capitalist production \textit{really} begins when production is extensive and a large number of workers are employed.

\item
(439) ``A large number of workers working together, at the same time, in one place (or, if you like, in the same field of labour), in order to produce the same sort of commodity under the command of the same capitalist, constitutes the starting point of capitalist production. This is true both historically and conceptually.''

\subitem
``At first, then, the difference is purely quantitative.''

\item
(443) ``Not only do we have here an increase in the productive power of the individual, by means of co-operation, but the creation of a new productive power, which is intrinsically a collective one.''

\subitem
(447) ``This power arises from co-operation itself. When the worker co-operates in a planned way with others, he strips off the fetters of his individuality, and develops the capabilities of his species.''

\item
(448--9) ``All directly social or communal labour on a large scale requires, to a greater or lesser degree, a directing authority.... A single violin player is his own conductor: an orchestra requires a separate one. The work of directing, superintending and adjusting becomes one of the functions of capital, from the moment that the labour under capital's control becomes co-operative. As a specific function of capital, the directing function requires its own special characteristics.''

\end{itemize}

\end{document}
